\section{Projektowanie}		%3

\hspace{0.60cm}Rozdział niniejszy opisuje proces projektowania architektury aplikacji mobilnej „Guess the Beat”, dobrane narzędzia programistyczne, kluczowe biblioteki języka Kotlin, obsługę sensorów urządzenia, integrację z usługami chmurowymi Firebase, schematy działania najważniejszych algorytmów odpowiedzialnych za przebieg rozgrywki oraz komunikację z zewnętrznym interfejsem programistycznym (API).

\subsection{Narzędzia programistyczne i kontrola wersji}

\hspace{0.60cm}Prace nad projektem realizowano z wykorzystaniem następującego oprogramowania:
\begin{itemize}
    \item \textbf{Android Studio:} Oficjalne zintegrowane środowisko programistyczne (IDE) dla systemu Android, zapewniające zaawansowane wsparcie dla języka Kotlin, wbudowany emulator (z możliwością symulacji czujników ruchu oraz czytnika linii papilarnych) oraz narzędzia do projektowania interfejsu w technologii Jetpack Compose.
    \item \textbf{Git oraz GitHub:} System kontroli wersji służący do śledzenia zmian w kodzie źródłowym, tworzenia gałęzi (\textit{branches}) oraz wspólnej pracy zespołowej nad projektem.
    \item \textbf{Konsola Firebase (Google Firebase Console):} Platforma chmurowa wykorzystywana do zarządzania kontami użytkowników, konfiguracji dostawców tożsamości (e-mail/hasło oraz Google OAuth 2.0) oraz administracji nierelacyjną bazą danych Cloud Firestore.
\end{itemize}

\subsection{Wybór bibliotek i architektura komunikacji wewnątrz aplikacji}

\hspace{0.60cm}Aplikacja została zaprojektowana w architekturze jednowidokowej (\textit{Single-Activity Architecture}) opartej na klasie \texttt{MainActivity} (dziedziczącej po \texttt{FragmentActivity} w celu wsparcia okien dialogowych biometrii), w której nawigacja między ekranami realizowana jest reaktywnie na podstawie stanu przechowywanego w komponencie \texttt{GameViewModel}. W projekcie wykorzystano następujące biblioteki i komponenty systemowe:

\begin{itemize}
    \item \textbf{Jetpack Compose (Material 3):} Nowoczesny, deklaratywny zestaw narzędzi do budowy natywnego interfejsu użytkownika. Zastosowano m.in. komponenty \texttt{Scaffold} z dolnym paskiem nawigacyjnym (\texttt{BottomNavBar}), \texttt{OutlinedTextField}, \texttt{LinearProgressIndicator}, \texttt{Card} oraz gradientowe tła (\texttt{Brush.verticalGradient}) wykorzystujące akcenty kolorystyczne motywu (\texttt{NeonBlue}).
    \item \textbf{Firebase Authentication oraz Cloud Firestore:} Zestaw bibliotek chmurowych Google służący do bezpiecznego uwierzytelniania użytkowników oraz trwałego przechowywania w chmurze ich indywidualnych statystyk (liczby rozegranych gier, odgadniętych i pominiętych utworów, sumy punktów oraz rekordu życiowego).
    \item \textbf{Android Credential Manager oraz Google Identity Services:} Nowoczesne API systemu Android (\texttt{androidx.credentials}) zapewniające natywną obsługę logowania za pomocą konta Google (\texttt{GetGoogleIdOption}) bez konieczności opuszczania ekranu aplikacji.
    \item \textbf{AndroidX Biometric (\texttt{BiometricPrompt}):} Biblioteka systemowa umożliwiająca weryfikację tożsamości użytkownika za pomocą czujników biometrycznych urządzenia (odcisk palca, skan twarzy) lub kodu blokady ekranu (\texttt{DEVICE\_CREDENTIAL}).
    \item \textbf{\texttt{MediaPlayer} oraz \texttt{Handler}:} Wbudowana klasa systemowa odpowiedzialna za strumieniowanie próbek audio z internetu (\texttt{prepareAsync}). Do precyzyjnego zatrzymywania dźwięku po upływie zadanego czasu etapu wykorzystano klasę \texttt{Handler} powiązaną z główną pętlą komunikatów (\texttt{Looper.getMainLooper()}) oraz metodą \texttt{postDelayed}.
    \item \textbf{\texttt{Kotlin Coroutines} (\texttt{Dispatchers.IO}) oraz \texttt{kotlinx-coroutines-tasks}:} Mechanizm współprogramów wykorzystywany do wykonywania zapytań sieciowych HTTP oraz asynchronicznych operacji Firebase (\texttt{.await()}) bez blokowania głównego wątku interfejsu użytkownika (\textit{UI Thread}).
    \item \textbf{Biblioteka Coil (\texttt{AsyncImage}):} Asynchroniczna biblioteka do ładowania i buforowania grafik z sieci, wykorzystywana do wyświetlania okładek albumów, playlist oraz zdjęć profilowych artystów.
    \item \textbf{\texttt{SensorManager} oraz \texttt{Sensor.TYPE\_ACCELEROMETER}:} Komponenty frameworka Android służące do nasłuchiwania zmian przyspieszenia urządzenia w trzech osiach ($X, Y, Z$) w celu wykrywania gestu odwrócenia telefonu.
\end{itemize}

\subsection{Projekt integracji z zewnętrznym API (Deezer)}

\hspace{0.60cm}Zamiast przechowywać pliki dźwiękowe lokalnie w pamięci urządzenia – co drastycznie zwiększyłoby rozmiar pakietu instalacyjnego APK i ograniczyło pulę utworów – aplikacja została zaprojektowana jako klient publicznego interfejsu \textbf{Deezer API}. Komunikacja sieciowa opiera się na wywołaniach protokołu HTTPS oraz kodowaniu zapytań w standardzie UTF-8 (\texttt{URLEncoder}).

\subsubsection{Przepływ danych i endpointy sieciowe}

\hspace{0.60cm}Klasa \texttt{MusicPlayer} komunikuje się z punktami końcowymi (ang. \textit{endpoints}) interfejsu Deezer API umożliwiającymi nie tylko wyszukiwanie playlist i utworów, ale również pobieranie aktualnych list przebojów (\textit{Charts}) dla topowych playlist, artystów oraz albumów, co przedstawia poniższy listing:

\begin{lstlisting}[caption={Wykorzystywane endpointy publicznego interfejsu Deezer API}, label={lst:deezer_endpoints}]
GET https://api.deezer.com/search/playlist?q={encodedQuery}
GET https://api.deezer.com/search?q={encodedQuery}&limit=5
GET https://api.deezer.com/playlist/{playlistId}
GET https://api.deezer.com/track/{trackId}?ts={timestamp}
GET https://api.deezer.com/chart/0/playlists
GET https://api.deezer.com/chart/0/artists
GET https://api.deezer.com/chart/0/albums
\end{lstlisting}

\hspace{0.60cm}Przedstawione na Listingu~\ref{lst:deezer_endpoints} zapytania HTTP realizują pełny cykl pobierania danych muzycznych w aplikacji. Pierwsze cztery endpointy odpowiadają za wyszukiwanie list odtwarzania, dynamiczne podpowiedzi utworów w trakcie wpisywania odpowiedzi, pobieranie identyfikatorów piosenek oraz metadanych konkretnego utworu. Trzy ostatnie punkty końcowe (\texttt{/chart/0/...}) służą do automatycznego zasilania menu głównego aktualnymi rankingami najpopularniejszych playlist, wykonawców i albumów.

\hspace{0.60cm}W celu ominięcia problemów z nieaktualną pamięcią podręczną (cache) po stronie serwera lub sieci, do zapytania o szczegóły konkretnego utworu dołączany jest dynamiczny parametr znacznika czasu (\texttt{ts=System.currentTimeMillis()}). Z odpowiedzi JSON dla obiektu \texttt{track} aplikacja wyodrębnia następujące pola:
\begin{itemize}
    \item \textbf{\texttt{title} oraz \texttt{artist.name}:} Tytuł utworu oraz nazwa wykonawcy, służące do weryfikacji odpowiedzi gracza oraz prezentacji wyników po zakończeniu rundy.
    \item \textbf{\texttt{preview}:} Bezpośredni adres URL do pliku MP3 zawierającego 30-sekundowy fragment utworu.
    \item \textbf{\texttt{album.cover\_medium}:} Adres URL okładki albumu o średniej rozdzielczości, przekazywany do komponentu \texttt{AsyncImage}.
\end{itemize}

\subsection{Projekt uwierzytelniania i bazy danych w chmurze (Firebase)}

\hspace{0.60cm}W celu umożliwienia graczom zapisywania swoich postępów między sesjami, aplikacja została zintegrowana z usługami \textbf{Firebase Authentication} oraz \textbf{Cloud Firestore}. Architektura warstwy danych zakłada trzy niezależne metody logowania:
\begin{itemize}
    \item \textbf{Klasyczne logowanie i rejestracja (E-mail i Hasło):} Realizowane za pomocą metod \texttt{signInWithEmailAndPassword} oraz \texttt{createUserWithEmailAndPassword} z walidacją minimalnej długości hasła ($6$ znaków).
    \item \textbf{Logowanie jednym kliknięciem przez konto Google:} Wykorzystuje nowoczesny mechanizm \texttt{CredentialManager} do pozyskania tokenu tożsamości (\texttt{GoogleIdTokenCredential}), który następnie wymieniany jest na poświadczenie sesji Firebase (\texttt{GoogleAuthProvider}).
    \item \textbf{Logowanie biometryczne (\texttt{BiometricAuthManager}):} Po pierwszym poprawnym zalogowaniu się adresem e-mail i hasłem, dane uwierzytelniające są zapisywane w prywatnej pamięci \texttt{SharedPreferences}. Przy kolejnych uruchomieniach aplikacji gracz może potwierdzić swoją tożsamość odciskiem palca lub skanem twarzy, co wyzwala automatyczne logowanie w tle.
\end{itemize}

\hspace{0.60cm}Statystyki każdego zarejestrowanego gracza przechowywane są w bazie dokumentowej Cloud Firestore w kolekcji \texttt{users}, pod dokumentem identyfikowanym przez unikalny identyfikator użytkownika (\texttt{uid}). Struktura dokumentu odwzorowuje klasę danych \texttt{PlayerStats} przechowującą pola: \texttt{gamesPlayed}, \texttt{songsGuessed}, \texttt{songsMissed}, \texttt{totalPoints} oraz \texttt{bestScore}.

\subsection{Opis kluczowych algorytmów aplikacji}

\hspace{0.60cm}Działanie rozgrywki opiera się na kilku dedykowanych algorytmach zaimplementowanych w warstwie logiki biznesowej:

\begin{enumerate}
    \item \textbf{Algorytm losowania puli rundowej i walidacji dostępności audio:} Po pobraniu listy identyfikatorów z wybranej playlisty, albumu lub dyskografii artysty, system weryfikuje, czy lista zawiera co najmniej $5$ utworów, a następnie losuje unikalną pulę pięciu piosenek na całą sesję (\texttt{ids.shuffled().take(totalRounds)}). Ze względu na ograniczenia licencyjne Deezer API, część utworów zwraca pusty ciąg znaków w polu \texttt{preview}. W takiej sytuacji algorytm automatycznie wybiera utwór zastępczy spoza bieżącej puli sesji (\texttt{availableReplacements.random()}) i rekurencyjnie ponawia ładowanie, gwarantując płynność rozgrywki bez zawieszania ekranu.
    
    \item \textbf{Algorytm progresywnego wydłużania czasu odsłuchu i punktacji malejącej:} Każda runda składa się z sześciu prób, którym przypisano rosnące przedziały czasowe odtwarzania próbki ($0,2\text{ s}$, $1\text{ s}$, $3\text{ s}$, $6\text{ s}$, $10\text{ s}$ oraz $15\text{ s}$) oraz malejącą pulę punktów możliwych do zdobycia (odpowiednio: $5000$, $4000$, $3000$, $2000$, $1000$ oraz $500$ punktów, co daje maksymalnie $25\,000$ punktów w pięciorundowej sesji). Każda błędna odpowiedź lub pominięcie próby wydłuża odsłuch kosztem niższej nagrody punktowej.
    
    \item \textbf{Algorytm detekcji obrotu telefonu z histerezą i blokadą czasową:} Moduł \texttt{FlipDetector} monitoruje składową przyspieszenia wzdłuż osi $Z$ akcelerometru. Wykrycie odwrócenia smartfona ekranem do dołu następuje po przekroczeniu progu $Z < -4,0\text{ m/s}^2$, pod warunkiem upływu minimum $1000\text{ ms}$ od poprzedniego zdarzenia. Ponowne wyzwolenie akcji jest blokowane flagą stanu (\texttt{isFaceDown}) aż do momentu przywrócenia telefonu do pozycji bazowej ($Z > 3,0\text{ m/s}^2$), co całkowicie eliminuje przypadkowe wielokrotne pominięcia etapów.
\end{enumerate}